% Generated draft: requires mathematical and source review.
\documentclass[11pt]{article}
\usepackage[margin=25mm]{geometry}
\usepackage{fontspec}
\setmainfont{Latin Modern Roman}
\usepackage{amsmath,amssymb,mathtools,unicode-math}
\setmathfont{Latin Modern Math}
\usepackage{enumitem,needspace,xurl,hyperref}
\setlength{\parindent}{0pt}
\setlength{\parskip}{5pt}
\setlength{\emergencystretch}{3em}
\newcommand{\N}{\mathbb{N}}
\newcommand{\Z}{\mathbb{Z}}
\newcommand{\Q}{\mathbb{Q}}
\newcommand{\R}{\mathbb{R}}
\newcommand{\C}{\mathbb{C}}
\newcommand{\F}{\mathbb{F}}
\newcommand{\A}{\mathbb{A}}
\newcommand{\PP}{\mathbb{P}}
\newcommand{\E}{\mathbb{E}}
\newcommand{\eps}{\varepsilon}
\newcommand{\dd}{\,\mathrm{d}}
\begin{document}
\section*{UnsolvedMath ID: 30000574}
\textbf{Source problem code:} catalogue entry
\subsection{Definitions and mathematical statement}
Let \(E\) be the nuclear Fr\'echet space in the source, over its intended scalar field \(\mathbb K\), usually \(\R\) or \(\C\). A Fr\'echet space is a complete metrizable locally convex Hausdorff vector space. One standard nuclearity formulation requires that for every continuous seminorm \(p\) there exist a continuous seminorm \(q\geq p\) for which the canonical map between the Banach completions \(E_q\to E_p\) is nuclear; these are completions of \(E/\ker q\) and \(E/\ker p\), respectively. The source does not give a narrower class of nuclear spaces.

For a continuous linear map \(S:E\to E\), a vector \(x\) is hypercyclic when \(\{S^nx:n\in\Z_{\geq0}\}\) is dense in \(E\), and \(S\) is hypercyclic if such a vector exists. The requested assertion, preserving the scalar quantifiers, is
\[
 \forall E\text{ in the intended class}\ \exists T\in\mathcal L(E)
 \ \forall\lambda\in\mathbb K\setminus\{0\}\ \exists x_\lambda\in E,
 \qquad
 \overline{\{\lambda^nT^nx_\lambda:n\geq0\}}=E.
\]
The same operator \(T\) must work for every nonzero scalar. The hypercyclic vector may depend on \(\lambda\); a common vector for all multiples is a stronger statement and is not requested. The case \(\lambda=1\) includes hypercyclicity of \(T\) itself. Scalars of arbitrary nonzero modulus are included, not merely scalars on the unit circle.

There is an essential space-scope omission: a dense orbit implies separability, and nonzero finite-dimensional real or complex spaces do not support hypercyclic linear operators in the usual topology, whereas finite-dimensional spaces are nuclear Fr\'echet spaces. The source says every nuclear Fr\'echet space without explicitly restricting to a separable infinite-dimensional setting. The literal broad quantifier therefore needs qualification; the draft records this issue rather than silently inserting the customary nontrivial hypotheses or treating the literal unrestricted version as a verified open conjecture. The treatment of the zero space and the scalar field are also unspecified.

The literature triage checked 2026-08-21 reports no authoritative resolution of the extracted claim located in its targeted search. The verification note says all quantifiers were preserved during notation cleanup; it does not supply the missing space restrictions. The open label remains dated source metadata, not an independently verified assertion about every space covered by the literal wording.
\subsection{Short English statement}
Seek one continuous linear operator whose every nonzero scalar multiple is hypercyclic; each multiple may have its own dense-orbit vector. The source omits separable infinite-dimensional restrictions, even though its literal nuclear Fr\'echet class includes finite-dimensional spaces. That scope issue is explicit rather than silently repaired.
\subsection{Sources}
\begin{itemize}
\item Local UnsolvedMath registry, permanent ID 30000574, source code catalogue entry, ``Scalar-Stable Hypercyclic Operators on Nuclear Fr\'echet Spaces,'' including the original and cleaned all-scalar question and quantifier-verification note.
\item \emph{Mini-Workshop: Hypercyclicity and Linear Chaos (2007)}, Oberwolfach Reports, pp.\ 2227--2276, DOI 10.4171/owr/2006/37: \url{https://doi.org/10.4171/owr/2006/37}. This is also the sole supplied literature reference. The citation uses 2007, while the DOI and proposed-year metadata use 2006.
\item Literature triage checked 2026-08-21; local record updated 2026-09-16. No external source has been visited or independently checked here.
\end{itemize}
\end{document}
